\subsection{商群}

定理1. 设R是带算子群G上的一个等价关系；若R与G上的群律左（相应地，右）相容（§3，第3号）且与Ω的作用相容，则e的等价类为G的一个稳定子群H，且关系R等价于 \(x^{-1}y \in H\) （相应地 \(yx^{-1} \in H\) ）。反之，若H为G的一个稳定子群，则关系 \(x^{-1}y \in H\) （相应地 \(yx^{-1} \in H\) 是一个等价关系，它在左侧（相应地，右侧）与 G 上的群运算相容，并且与 Ω 的作用相容，在此关系下 H 是 e 的等价类。

我们将注意力限制在关系 R 与 G 上的运算左相容的情形（右相容关系的情形可通过将 G 上的运算替换为相反运算来得到）。该关系 \(y \equiv x \pmod{R}\) 等价于 \(x^{-1}y \equiv e \pmod{R}\) ，对于 \(y \equiv x\) 蕴含 \(x^{-1}y \equiv x^{-1}x = e\) 反之亦然 \(x^{-1}y \equiv e\) 蕴含 \(y = x(x^{-1}y) \equiv x\) 。若 H 表示 e 的等价类，则关系 R 等价于 \(x^{-1}y \in H\) 。我们证明 H 是 G 的一个稳定子群。对于每个算子 \(\alpha\) ，关系 \(x \equiv e\) 蕴含 \(x^{\alpha} \equiv e^{\alpha} = e\) ，因此 \(H^{\alpha} \subset H\) 且 H 在 \(\Omega\) 的作用下是稳定的。只需建立（命题 1） \(x \in H\) 和 \(y \in H\) 蕴含 \(x^{-1}y \in H\) ，即 \(x \equiv e\) 和 \(y \equiv e\) 蕴含 \(x \equiv y\) ，这是 R 的传递性的一个推论。

反之，设H是G的一个稳定子群；该关系 \(x^{-1}y \in H\) 是自反的，因为 \(x^{-1}x = e \in H\) ；它是对称的，因为 \(x^{-1}y \in H\) 蕴含 \(y^{-1}x = (x^{-1}y)^{-1} \in H\) 它是传递的，因为 \(x^{-1}y \in H\) 和 \(y^{-1}z \in H\) 蕴含 \(x^{-1}z = (x^{-1}y)(y^{-1}z) \in H\) 它与G上的合成律左相容，因为 \(x^{-1}y = (zx)^{-1}(zy)\) 为所有 \(z \in G\) 最后，对于每个算子 \(\alpha\) ，关系 \(y \in xH\) 蕴含 \(y^{\alpha} \in x^{\alpha}H^{\alpha} \subset x^{\alpha}H\) 因此等价关系 \(x^{-1}y \in H\) 与 \(\Omega\) 在G上的作用相容。

设G是一个群，H是G的一个子群；关系 \(x^{-1}y \in H\) （相应地 \(yx^{-1} \in H\) )也写成等价形式 \(y \in xH\) （相应地 \(y \in Hx\) ). 因此，G 的每个子群 H 都在 G 上定义了两种等价关系，即 \(y \in xH\) 和 \(y \in Hx\) 在这些关系下的等价类分别是集合xH，称为H的左陪集（或模H的左陪集），以及集合Hx，称为H的右陪集（或模H的右陪集）。通过关于这些关系对子集A ⊂ G进行饱和（集合论，II，§6，第4节），我们分别得到集合AH和HA。映射 \(x \mapsto x^{-1}\) 将模H的左陪集变换为模H的右陪集，反之亦然。

左陪集（模 H）的集合的基数称为子群 H 关于 G 的指数，记为 (G:H)；它也等于右陪集的集合的基数。

若G的子群K包含H，则K是H的左（或右）陪集的并集。由于K的左陪集由K经左平移得到，包含在K的一个左陪集内的H的左陪集所成集合的基数与后者无关。因此（集合论，III，§5，第8节，命题9）：

命题4。设H和K是群G的两个子群，且H ⊂ K。则

\[
(\mathbf {G}: \mathbf {H}) = (\mathbf {G}: \mathbf {K}) (\mathbf {K}: \mathbf {H}).\tag{2}
\]

推论。若G是一个阶为g的有限群，且H是G的一个阶为h的子群，则

\[
h. (\mathbf {G}: \mathbf {H}) = g\tag{3}
\]

（特别地，H 的阶和指数都是 G 的阶的因数。）

定理1使我们能够确定与带算子群G上的运算相容的等价关系：若R是这样一个关系，则它既与G上的群运算左相容又右相容，并且与算子的作用相容。 \(\Omega\) 因此，若H是e（模R）的类，则H是一个稳定子群，使得关系 \(y \in xH\) 和 \(y \in Hx\) 是等价的（因为两者都等价于R）；因此 \(xH = Hx\) 为所有 \(x \in G\) 。反之，若如此，则等价关系之一 \(y \in xH\) , \(y \in Hx\) 与群律相容，因为它与该律既左相容又右相容（§3，第4节），并且与 \(\Omega\) 的作用相容。由于方程 \(xH = Hx\) 等价于 \(xHx^{-1} = H\) ，我们作出如下定义：

定义5. 设G是一个带算子的群。G的一个稳定子群H称为G的一个正规（或不变）稳定子群，如果 \(xHx^{-1} = H\) 为所有 \(x \in G\) .

上定义了一个处处有定义的合成律。若 \(\Omega = \varnothing\) ，G的一个正规稳定子群称为G的一个正规（或不变）子群。在交换群中，每个子群都是正规的。

要验证一个稳定子群H是正规的，只需证明 \(xHx^{-1} \subset H\) 为所有 \(x \in G\) ；因为若如此，则 \(x^{-1}Hx \subset H\) 为所有 \(x \in G\) ，即 \(H \subset xHx^{-1}\) ，从而 \(H = xHx^{-1}\) .

设H是G的一个正规稳定子群，R是由H定义的等价关系 \(y \in xH\) ；在商集G/R上，内运算（即群G的运算关于R的商）是结合的；e的类是此商运算的单位元；G中两个互逆元素的类在商运算下互为逆元，并且 \(\Omega\) 的作用（即 \(\Omega\) 在G上的作用关于R的商）关于G/R上的内运算是分配的（ \(\S 3\) ，第5节）。因此，总结所得结果：

定理2. 设G是一个带算子的群。对于G上的一个等价关系R，要使它与群运算及 \(\Omega\) 的作用相容，必要且充分条件是它具有形式 \(x^{-1}y \in H\) ，其中H是G的一个正规稳定子群（对于这样的子群，关系 \(x^{-1}y \in H\) 还等价于 \(yx^{-1} \in H\) ）。G/R上的合成运算（即G上运算的商）以及 \(\Omega\) 在G/R上的作用（即 \(\Omega\) 在G上的作用关于这样的关系R的商）赋予G/R带算子的群的结构，称为商结构，并且到商集的典范映射是带算子的群的同态。

定义6. 带算子的群G关于由G的正规子群H定义的等价关系的商，连同商结构，称为G关于H的带算子的商群，记为G/H。典范映射G → G/H称为典范同态。

设G是一个群，H是G的一个正规子群。商G/H连同其群结构称为G关于H的商群。对于从G/H到某个带算子的群的映射，要使它是带算子的群的同态，必要且充分条件是它与从G到G/H的典范映射的复合是同态：这证明了“商群”这一名称的合理性（集合论，IV，§ 2，第6节）。

由G的正规稳定子群定义的等价关系记为 \(x \equiv y \pmod{H}\) 或 \(x \equiv y(H)\) .

命题5. 设 \(f: G \to G'\) 是带算子的群的同态，H和H'分别是G和G'的正规稳定子群，使得 \(f(H) \subset H'\) 该映射 \(f\) 与由H和H'定义的等价关系相容。设 \(\pi: G \to G/H\) 和 \(\pi': G' \to G'/H'\) 是典范同态。由 \(\bar{f}: G/H \to G'/H'\) 通过取商导出的映射 \(f\) 是一个同态。

上定义了一个处处有定义的合成律。若 \(x \equiv y \pmod{H}\) ，则 \(x^{-1}y \in H\) 的自同态，因此

\[
f (x) ^ {- 1} f (y) = f \left(x ^ {- 1}\right) f (y) = f \left(x ^ {- 1} y\right) \in f (\mathrm{H}) \subset \mathrm{H} ^ {\prime}
\]

，因此 \(f(x) \equiv f(y) \pmod{H'}\) 第二个断言由商律的泛性质（§1，第6节）得出。

注记。(1) 若A是群G的任意子集，H是G的正规子群，则AH = HA；该集合通过对A关于该关系进行饱和而得到。 \(x \equiv y \pmod{H}\) .

\begin{enumerate}
\def\labelenumi{(\arabic{enumi})}
\setcounter{enumi}{1}
\tightlist
\item
  若H是G的有限指数的正规子群，则商群G/H是阶为(G:H)的有限群。
\end{enumerate}

注意，若H是群G的正规子群，且K是H的正规子群，则K不一定是G的正规子群（I，§5，习题10）。

设G是一个带算子的群。G的所有正规稳定子群的任意族的交是正规稳定子群。因此，对于G的每个子集X，存在包含X的最小正规稳定子群，称为由X生成的正规稳定子群。

在具有算子群 G 中，稳定子群 G 和 \(\{e\}\) 是正规的。

定义 7. 具有算子群 G 称为单的，如果 \(G \neq \{e\}\) 并且 G 中不存在除 G 和 \(\{e\}\) .

